% 実験方法は過去形で書くこと（重要）
% サブセクションを使って操作手順を構造化する
% 番号付きリスト(enumerate)を使用して手順を明確にする
% 【例】回路作成の場合:
% \subsection{回路の作成}
% \subsubsection{実験器具}
%   \begin{enumerate}
%     \item 測定機器\\
%       \quad 定電圧直流電源、ファンクションジェネレータ、オシロスコープ、テスタ
%     \item 回路部品\\
%       \quad NPN トランジスタ(2CS1815)、\SI{10}{\micro F} 電解コンデンサ、
%       抵抗器(\SI{1}{\kilo\ohm}、\SI{10}{\kilo\ohm})
%   \end{enumerate}
% \subsubsection{回路作成手順}
%   \begin{enumerate}
%     \item 実体配線図の作成と部品の取り付け\\
%       \quad ユニバーサル基板の表面と裏面に部品配置を鉛筆で描いた。
%       実体配線図に基づいて、部品をユニバーサル基板に差し込み、ハンダ付けを行った。
%     \item 導通チェック\\
%       \quad テスタで導通を確認し、GND 間の導通、抵抗端子の導通が正しいことを確認した。
%   \end{enumerate}

\section{実験方法}
% ここに実験手順を過去形で詳細に記載してください。
% 他者が再現できるよう、条件や数値を明確に記述します。

\subsection{使用機器・材料}
% 使用した機器や材料をリスト形式で記載します。
  \begin{enumerate}
    \item 測定機器\\
      \quad % ここに測定機器を列挙
    \item 試薬・材料\\
      \quad % ここに試薬や材料を列挙
  \end{enumerate}

\subsection{実験手順}
% 実験の手順を過去形で記載します。
  \begin{enumerate}
    \item % 手順1を記載
    \item % 手順2を記載
    \item % 手順3を記載
  \end{enumerate}
